\documentclass[../../main.tex]{subfiles}
\graphicspath{{\subfix{../../image/}}}

\begin{document}

\section{不动点方法}\label{section:ch09-2-不动点方法}

接着研究拓扑不动点定理对非线性偏微分方程的适用性。有用的抽象定理至少有三类：
\begin{enumerate}[(a)]
\item 严格压缩映射的不动点定理；
\item 紧映射的不动点定理；
\item 保序算子的不动点定理。
\end{enumerate}
下面给出类型 (a) 与 (b) 的应用。保序性质对非线性 PDE 的用处将在 \S9.3 说明。

\subsection{Banach 不动点定理}\label{subsection:ch09-2-1-Banach}

此后 $X$ 表示一个 Banach 空间。最简单的不动点定理是

\begin{theorem}[Banach 不动点定理]\label{theorem:ch09-2-1-Banach}
设 $A:X\to X$ 为非线性映射，且假设
\begin{align}
\|A[u]-A[\widetilde{u}]\|\leqslant\gamma\|u-\widetilde{u}\|\qquad (u,\widetilde{u}\in X)\label{align:ch09-2-eq01}
\end{align}
对某常数 $\gamma<1$。则 $A$ 有唯一不动点。
\begin{definition}
若\eqref{align:ch09-2-eq01} 成立，则称 $A$ 为\textbf{严格压缩映射}。
\end{definition}
\begin{proof}
固定任一点 $u_0\in X$，随后迭代定义 $u_{k+1}=A[u_k]$，$k=0,1,\cdots$。则
\begin{align*}
\|A[u_{k+1}]-A[u_k]\|\leqslant\gamma\|u_{k+1}-u_k\|=\gamma\|A[u_k]-A[u_{k-1}]\|,
\end{align*}
故
\begin{align*}
\|A[u_{k+1}]-A[u_k]\|\leqslant\gamma^k\|A[u_0]-u_0\|\qquad (k=1,2,\cdots).
\end{align*}
于是若 $k\geqslant l$，
\begin{align*}
\|u_k-u_l\|=\|A[u_{k-1}]-A[u_{l-1}]\|\leqslant\sum_{j=l-1}^{k-2}\|A[u_{j+1}]-A[u_j]\|\leqslant\|A[u_0]-u_0\|\sum_{j=l-1}^{k-2}\gamma^j.
\end{align*}
故 $\{u_k\}_{k=1}^\infty$ 是 $X$ 中的 Cauchy 列，从而存在点 $u\in X$ 使 $u_k\to u$ 于 $X$。显然 $A[u]=u$。故 $u$ 是 $A$ 的不动点，且假设\eqref{align:ch09-2-eq01} 保证唯一性。
\end{proof}
\end{theorem}

Banach 不动点定理对非线性 PDE 的应用通常涉及各种摄动论证：给定一个性态良好的线性椭圆偏微分方程，把一个小的非线性修正写成压缩映射往往是直接的。此类证明的标志是出现一个必须取充分小以保证严格压缩性质的参数。

然而有时可把给定问题分解为若干子问题，每个子问题带一个小参数：

\begin{example}[反应扩散方程]\label{example:ch09-2-1-反应扩散}
考察反应扩散方程组初边值问题的可解性：
\begin{align}
\begin{cases}
\mathbf{u}_t-\Delta\mathbf{u}=\mathbf{f}(\mathbf{u})&\text{在 }U_T\text{ 中},\\
\mathbf{u}=0&\text{在 }\partial U\times[0,T]\text{ 上},\\
\mathbf{u}=\mathbf{g}&\text{在 }U\times\{t=0\}\text{ 上}.
\end{cases}\label{align:ch09-2-eq02}
\end{align}
此处 $\mathbf{u}=(u_1,\cdots,u_m)$，$\mathbf{g}=(g_1,\cdots,g_m)$，如常 $U_T=U\times(0,T]$，$U\subset\mathbb{R}^n$ 为开、有界、边界光滑区域。时间 $T>0$ 固定。假设初值函数 $\mathbf{g}\in H_0^1(U;\mathbb{R}^m)$。关于非线性项，设
\begin{align}
\mathbf{f}:\mathbb{R}^m\to\mathbb{R}^m\text{ 是 Lipschitz 连续的}.\label{align:ch09-2-eq03}
\end{align}
此假设特别蕴涵
\begin{align}
|\mathbf{f}(z)|\leqslant C(1+|z|)\label{align:ch09-2-eq04}
\end{align}
对每个 $z\in\mathbb{R}^m$ 及某常数 $C$。

采用 \S7.1 的术语，称函数
\begin{align}
\mathbf{u}\in L^2(0,T;H_0^1(U;\mathbb{R}^m)),\quad \mathbf{u}'\in L^2(0,T;H^{-1}(U;\mathbb{R}^m))\label{align:ch09-2-eq05}
\end{align}
为\eqref{align:ch09-2-eq02} 的弱解，只要
\begin{align}
\langle\mathbf{u}',\mathbf{v}\rangle+B[\mathbf{u},\mathbf{v}]=(\mathbf{f}(\mathbf{u}),\mathbf{v})\quad\text{a.e. }0\leqslant t\leqslant T\label{align:ch09-2-eq06}
\end{align}
对每个 $\mathbf{v}\in H_0^1(U;\mathbb{R}^m)$ 成立，且
\begin{align}
\mathbf{u}(0)=\mathbf{g}.\label{align:ch09-2-eq07}
\end{align}
在\eqref{align:ch09-2-eq06} 中 $\langle,\rangle$ 表示 $H^{-1}(U;\mathbb{R}^m)$ 与 $H_0^1(U;\mathbb{R}^m)$ 的配对，$B[,]$ 是 $-\Delta$ 在 $H_0^1(U;\mathbb{R}^m)$ 中对应的双线性型，$(,)$ 表示 $L^2(U;\mathbb{R}^m)$ 的内积。$H_0^1(U;\mathbb{R}^m)$ 中的范数取为
\begin{align*}
\|\mathbf{u}\|_{H_0^1(U;\mathbb{R}^m)}=\left(\int_U|D\mathbf{u}|^2\,\mathrm{d}x\right)^{1/2}.
\end{align*}
由 \S5.9.2 回顾，\eqref{align:ch09-2-eq05} 蕴涵在一个零测集上重新定义 $\mathbf{u}$ 后，$\mathbf{u}\in C([0,T];L^2(U;\mathbb{R}^m))$。

\begin{theorem}[存在性]\label{theorem:ch09-2-2-反应扩散存在}
\eqref{align:ch09-2-eq02} 存在唯一弱解。
\begin{proof}
1. 将在空间 $X=C([0,T];L^2(U;\mathbb{R}^m))$ 中应用 Banach 定理，范数为
\begin{align*}
\|\mathbf{v}\|\triangleq\max_{0\leqslant t\leqslant T}\|\mathbf{v}(t)\|_{L^2(U;\mathbb{R}^m)}.
\end{align*}
算子 $A$ 定义如下。给定 $\mathbf{u}\in X$，令 $\mathbf{h}(t)=\mathbf{f}(\mathbf{u}(t))$（$0\leqslant t\leqslant T$）。由增长估计\eqref{align:ch09-2-eq04}，$\mathbf{h}\in L^2(0,T;L^2(U;\mathbb{R}^m))$。故 \S7.1 的理论保证线性抛物 PDE
\begin{align}
\begin{cases}
\mathbf{w}_t-\Delta\mathbf{w}=\mathbf{h}&\text{在 }U_T\text{ 中},\\
\mathbf{w}=0&\text{在 }\partial U\times[0,T]\text{ 上},\\
\mathbf{w}=\mathbf{g}&\text{在 }U\times\{t=0\}\text{ 上}
\end{cases}\label{align:ch09-2-eq08}
\end{align}
有唯一弱解
\begin{align}
\mathbf{w}\in L^2(0,T;H_0^1(U;\mathbb{R}^m)),\quad \mathbf{w}'\in L^2(0,T;H^{-1}(U;\mathbb{R}^m)).\label{align:ch09-2-eq09}
\end{align}
于是 $\mathbf{w}\in X$ 满足
\begin{align}
\langle\mathbf{w}',\mathbf{v}\rangle+B[\mathbf{w},\mathbf{v}]=(\mathbf{h},\mathbf{v})\quad\text{a.e. }0\leqslant t\leqslant T\label{align:ch09-2-eq10}
\end{align}
对每个 $\mathbf{v}\in H_0^1(U;\mathbb{R}^m)$ 成立，且 $\mathbf{w}(0)=\mathbf{g}$。定义 $A:X\to X$ 为 $A[\mathbf{u}]=\mathbf{w}$。
2. 现断言
\begin{align}
\text{若 }T>0\text{ 充分小，则 }A\text{ 是严格压缩映射。}\label{align:ch09-2-eq11}
\end{align}
为证此，取 $\mathbf{u},\widetilde{\mathbf{u}}\in X$，如上定义 $\mathbf{w}=A[\mathbf{u}]$、$\widetilde{\mathbf{w}}=A[\widetilde{\mathbf{u}}]$。于是 $\mathbf{w}$ 对 $\mathbf{h}=\mathbf{f}(\mathbf{u})$ 满足\eqref{align:ch09-2-eq10}，而 $\widetilde{\mathbf{w}}$ 对 $\widetilde{\mathbf{h}}\triangleq\mathbf{f}(\widetilde{\mathbf{u}})$ 满足类似恒等式。如 \S7.1 计算，
\begin{align}
\frac{\mathrm{d}}{\mathrm{d}t}\|\mathbf{w}-\widetilde{\mathbf{w}}\|_{L^2}^2+2\|\mathbf{w}-\widetilde{\mathbf{w}}\|_{H_0^1}^2&=2(\mathbf{w}-\widetilde{\mathbf{w}},\mathbf{h}-\widetilde{\mathbf{h}})\notag\\
&\leqslant\varepsilon\|\mathbf{w}-\widetilde{\mathbf{w}}\|_{L^2}^2+\frac1\varepsilon\|\mathbf{f}(\mathbf{u})-\mathbf{f}(\widetilde{\mathbf{u}})\|_{L^2}^2\notag\\
&\leqslant\varepsilon C\|\mathbf{w}-\widetilde{\mathbf{w}}\|_{H_0^1}^2+\frac1\varepsilon\|\mathbf{f}(\mathbf{u})-\mathbf{f}(\widetilde{\mathbf{u}})\|_{L^2}^2,\label{align:ch09-2-eq12}
\end{align}
最后一步由 Poincaré 不等式成立。取 $\varepsilon>0$ 充分小，推得
\begin{align*}
\frac{\mathrm{d}}{\mathrm{d}t}\|\mathbf{w}-\widetilde{\mathbf{w}}\|_{L^2}^2\leqslant C\|\mathbf{f}(\mathbf{u})-\mathbf{f}(\widetilde{\mathbf{u}})\|_{L^2}^2\leqslant C\|\mathbf{u}-\widetilde{\mathbf{u}}\|_{L^2}^2,
\end{align*}
因 $\mathbf{f}$ Lipschitz。故
\begin{align}
\|\mathbf{w}(s)-\widetilde{\mathbf{w}}(s)\|_{L^2}^2\leqslant C\int_0^s\|\mathbf{u}(t)-\widetilde{\mathbf{u}}(t)\|_{L^2}^2\,\mathrm{d}t\leqslant CT\|\mathbf{u}-\widetilde{\mathbf{u}}\|^2\label{align:ch09-2-eq13}
\end{align}
对每个 $0\leqslant s\leqslant T$。关于 $s$ 最大化左端，得
\begin{align*}
\|\mathbf{w}-\widetilde{\mathbf{w}}\|^2\leqslant CT\|\mathbf{u}-\widetilde{\mathbf{u}}\|^2.
\end{align*}
故
\begin{align}
\|A[\mathbf{u}]-A[\widetilde{\mathbf{u}}]\|\leqslant(CT)^{1/2}\|\mathbf{u}-\widetilde{\mathbf{u}}\|,\label{align:ch09-2-eq14}
\end{align}
于是若 $T>0$ 充分小使 $(CT)^{1/2}=\gamma<1$，$A$ 是严格压缩映射。
3. 给定任一 $T>0$，取 $T_1>0$ 充分小使 $(CT_1)^{1/2}<1$。然后可应用 Banach 不动点定理，找到在时间区间 $[0,T_1]$ 上存在的\eqref{align:ch09-2-eq02} 的弱解 $\mathbf{u}$。因 $\mathbf{u}(t)\in H_0^1(U;\mathbb{R}^m)$ 对 a.e. $0\leqslant t\leqslant T_1$，必要时可重新定义 $T_1$，设 $\mathbf{u}(T_1)\in H_0^1(U;\mathbb{R}^m)$。注意时间 $T_1>0$ 仅依赖于 $\mathbf{f}$ 的 Lipschitz 常数。故可重复上述论证，把解延拓到时间区间 $[T_1,2T_1]$。继续，经有限步后构造出在整个区间 $[0,T]$ 上存在的弱解。
4. 为证唯一性，设 $\mathbf{u}$ 与 $\widetilde{\mathbf{u}}$ 均为\eqref{align:ch09-2-eq02} 的弱解。则在不等式\eqref{align:ch09-2-eq13} 中有 $\mathbf{w}=\mathbf{u}$、$\widetilde{\mathbf{w}}=\widetilde{\mathbf{u}}$，从而
\begin{align*}
\|\mathbf{u}(s)-\widetilde{\mathbf{u}}(s)\|_{L^2}^2\leqslant C\int_0^s\|\mathbf{u}(t)-\widetilde{\mathbf{u}}(t)\|_{L^2}^2\,\mathrm{d}t
\end{align*}
对 $0\leqslant s\leqslant T$。由 Gronwall 不等式，$\mathbf{u}\equiv\widetilde{\mathbf{u}}$。
\end{proof}
\end{theorem}
\end{example}

\paragraph{解释。}在常见应用中，问题\eqref{align:ch09-2-eq02} 记录各种化学物质浓度 $u_1,\cdots,u_m$ 的演化，它们既在介质中扩散又相互反应。扩散项为 $\Delta\mathbf{u}$（更一般地为 $(a_1\Delta u_1,\cdots,a_m\Delta u_m)$，其中常数 $a_k>0$ 刻画第 $k$ 种化学物质的扩散）。反应项 $\mathbf{f}(\mathbf{u})$ 建立化学模型。在上述例子中我们作了 $\mathbf{f}$ 全局 Lipschitz 这一不切实际的假设。在更现实的模型中，$\mathbf{f}$ 常为 $\mathbf{u}$ 的多项式，关于解的整体存在或爆破存在有趣问题。（一个简单例子在 \S9.4.1 处理。）

后面将在 \S12.2 用 Banach 不动点定理证明某些拟线性波方程解的长时与短时存在性。

\subsection{Schauder 不动点定理、Schaefer 不动点定理}\label{subsection:ch09-2-2-SchauderSchaefer}

接着把 Brouwer 不动点定理（\S8.1.4）推广到 Banach 空间。关键假设现在是紧性。本节中 $X$ 继续表示一个实 Banach 空间。

\begin{theorem}[Schauder 不动点定理]\label{theorem:ch09-2-3-Schauder}
设 $K\subset X$ 非空、紧、凸，且 $A:K\to K$ 连续。则 $A$ 在 $K$ 中有不动点。
\begin{proof}
1. 固定 $\varepsilon>0$，选取有限个点 $u_1,\cdots,u_{N_\varepsilon}\in K$，使开球 $\{B_0(u_i,\varepsilon)\}_{i=1}^{N_\varepsilon}$ 覆盖 $K$：
\begin{align}
K\subset\bigcup_{i=1}^{N_\varepsilon}B_0(u_i,\varepsilon).\label{align:ch09-2-eq15}
\end{align}
这是可能的，因 $K$ 紧。令 $K_\varepsilon$ 表示点 $\{u_1,\cdots,u_{N_\varepsilon}\}$ 的闭凸包：
\begin{align*}
K_\varepsilon\triangleq\left\{\sum_{i=1}^{N_\varepsilon}\lambda_i u_i\mid 0\leqslant\lambda_i\leqslant 1,\ \sum_{i=1}^{N_\varepsilon}\lambda_i=1\right\}.
\end{align*}
因 $K$ 凸，$K_\varepsilon\subset K$。现定义 $P_\varepsilon:K\to K_\varepsilon$ 为
\begin{align*}
P_\varepsilon[u]\triangleq\frac{\sum_{i=1}^{N_\varepsilon}\operatorname{dist}(u,K\setminus B_0(u_i,\varepsilon))u_i}{\sum_{i=1}^{N_\varepsilon}\operatorname{dist}(u,K\setminus B_0(u_i,\varepsilon))}\qquad (u\in K).
\end{align*}
由\eqref{align:ch09-2-eq15}，分母永不为零。显然 $P_\varepsilon$ 连续，且对每个 $u\in K$，
\begin{align}
\|P_\varepsilon[u]-u\|\leqslant\frac{\sum_{i=1}^{N_\varepsilon}\operatorname{dist}(u,K\setminus B_0(u_i,\varepsilon))\|u_i-u\|}{\sum_{i=1}^{N_\varepsilon}\operatorname{dist}(u,K\setminus B_0(u_i,\varepsilon))}\leqslant\varepsilon.\label{align:ch09-2-eq16}
\end{align}
2. 接着考虑算子 $A_\varepsilon:K_\varepsilon\to K_\varepsilon$，定义为
\begin{align*}
A_\varepsilon[u]\triangleq P_\varepsilon[A[u]]\qquad (u\in K_\varepsilon).
\end{align*}
$K_\varepsilon$ 同胚于 $\mathbb{R}^{M_\varepsilon}$ 中的闭单位球，某 $M_\varepsilon\leqslant N_\varepsilon$。故 Brouwer 不动点定理（\S8.1.4）保证存在点 $u_\varepsilon\in K_\varepsilon$ 使
\begin{align*}
A_\varepsilon[u_\varepsilon]=u_\varepsilon.
\end{align*}
3. 因 $K$ 紧，存在子列 $\varepsilon_j\to 0$ 及点 $u\in K$，使 $u_{\varepsilon_j}\to u$ 于 $X$。断言 $u$ 是 $A$ 的不动点。事实上，用估计\eqref{align:ch09-2-eq16}，推得
\begin{align*}
\|u_{\varepsilon_j}-A[u_{\varepsilon_j}]\|=\|A_{\varepsilon_j}[u_{\varepsilon_j}]-A[u_{\varepsilon_j}]\|=\|P_{\varepsilon_j}[A[u_{\varepsilon_j}]]-A[u_{\varepsilon_j}]\|\leqslant\varepsilon_j.
\end{align*}
于是因 $A$ 连续，得 $u=A[u]$。
\end{proof}
\end{theorem}

接着把 Schauder 不动点定理改写成另一种形式，后者对非线性偏微分方程应用常常更有用。

\begin{definition}\label{definition:ch09-2-紧映射}
称非线性映射 $A:X\to X$ 是\textbf{紧的}，只要对每个有界序列 $\{u_k\}_{k=1}^\infty$，序列 $\{A[u_k]\}_{k=1}^\infty$ 是预紧的；即存在子列 $\{u_{k_j}\}_{j=1}^\infty$ 使 $\{A[u_{k_j}]\}_{j=1}^\infty$ 在 $X$ 中收敛。
\end{definition}

\begin{theorem}[Schaefer 不动点定理]\label{theorem:ch09-2-4-Schaefer}
设 $A:X\to X$ 连续且紧。又设集合
\begin{align*}
\{u\in X\mid u=\lambda A[u]\text{ 对某 }0\leqslant\lambda\leqslant 1\}
\end{align*}
有界。则 $A$ 有不动点。
\begin{proof}
断言是说：若对所有算子 $\lambda A$（$0\leqslant\lambda\leqslant 1$）的任何可能不动点有一致界，则 $A$ 有不动点存在。这符合那个著名的非正式原理："若能在假设非线性 PDE 解存在的条件下证明关于这些解的适当估计，则事实上这些解确实存在"。这就是\textbf{先验估计}方法。
Schaefer 定理对 Schauder 定理的应用优势在于：不必识别一个显式的凸、紧集合。
1. 选常数 $M$ 充分大使
\begin{align}
\|u\|<M\quad\text{若 }u=\lambda A[u]\text{ 对某 }0\leqslant\lambda\leqslant 1.\label{align:ch09-2-eq17}
\end{align}
再定义
\begin{align}
\widetilde{A}[u]\triangleq\begin{cases}A[u]&\text{若 }\|A[u]\|\leqslant M,\\ \dfrac{MA[u]}{\|A[u]\|}&\text{若 }\|A[u]\|\geqslant M.\end{cases}\label{align:ch09-2-eq18}
\end{align}
注意 $\widetilde{A}:B(0,M)\to B(0,M)$。令 $K=\widetilde{A}(B(0,M))$ 的闭凸包。则因 $A$ 从而 $\widetilde{A}$ 是紧映射，$K$ 是 $X$ 的紧凸子集。且 $\widetilde{A}:K\to K$。
2. 调用 Schauder 不动点定理，推得存在点 $u\in K$ 使
\begin{align}
\widetilde{A}[u]=u.\label{align:ch09-2-eq19}
\end{align}
现还断言 $u$ 是 $A$ 的不动点。否则由\eqref{align:ch09-2-eq18} 与\eqref{align:ch09-2-eq19}，将有
\begin{align*}
\|A[u]\|>M
\end{align*}
且
\begin{align}
u=\lambda A[u]\quad\text{其中 }\lambda=\frac{M}{\|A[u]\|}<1.\label{align:ch09-2-eq20}
\end{align}
但 $\|u\|=\|\widetilde{A}[u]\|=M$，与\eqref{align:ch09-2-eq17} 及\eqref{align:ch09-2-eq20} 矛盾。
\end{proof}
\end{theorem}

\paragraph{凸集上的不动点定理。}对某些应用，下面这个前述定理的变体是方便的：若 $K$ 是 Banach 空间 $X$ 的凸子集，$0\in K$，且 $A:K\to K$ 是连续紧映射，集合
\begin{align*}
\{u\in K\mid u=\lambda A[u]\text{ 对某 }0\leqslant\lambda\leqslant 1\}
\end{align*}
有界，则 $A$ 在 $K$ 中有不动点。证明是上述论证的一个小修改。

\paragraph{应用。}用 Schauder 与 Schaefer 不动点定理解 PDE 依赖于与 Banach 定理应用相当不同的考虑。关键假设现在不是某个参数小，而是我们有某种紧性。因线性椭圆算子的逆通常是磨光的，对某些非线性椭圆方程确实可得到紧性。下面是一个快速但相当粗糙的例子：

\begin{example}[一个拟线性椭圆 PDE]\label{example:ch09-2-2-拟线性椭圆}
现给出 Schaefer 定理的一个简单应用，求解半线性边值问题
\begin{align}
\begin{cases}
-\Delta u+b(Du)+\mu u=0&\text{在 }U\text{ 中},\\
u=0&\text{在 }\partial U\text{ 上},
\end{cases}\label{align:ch09-2-eq21}
\end{align}
其中 $U$ 有界、$\partial U$ 光滑。设 $b:\mathbb{R}^n\to\mathbb{R}$ 光滑、Lipschitz 连续，从而满足增长条件
\begin{align}
|b(p)|\leqslant C(|p|+1)\label{align:ch09-2-eq22}
\end{align}
对某常数 $C$ 及所有 $p\in\mathbb{R}^n$。

\begin{theorem}[存在性]\label{theorem:ch09-2-5-半线性椭圆存在}
若 $\mu>0$ 充分大，则存在函数 $u\in H^2(U)\cap H_0^1(U)$ 求解边值问题\eqref{align:ch09-2-eq21}。
\begin{proof}
1. 给定 $u\in H_0^1(U)$，令
\begin{align}
f\triangleq-b(Du).\label{align:ch09-2-eq23}
\end{align}
由估计\eqref{align:ch09-2-eq22}，$f\in L^2(U)$。现令 $w\in H_0^1(U)$ 为线性问题
\begin{align}
\begin{cases}
-\Delta w+\mu w=f&\text{在 }U\text{ 中},\\
w=0&\text{在 }\partial U\text{ 上}
\end{cases}\label{align:ch09-2-eq24}
\end{align}
的唯一弱解。由 \S6.3 证明的正则性理论，还知道 $w\in H^2(U)$，且估计
\begin{align}
\|w\|_{H^2(U)}\leqslant C\|f\|_{L^2(U)}\label{align:ch09-2-eq25}
\end{align}
对某常数 $C$。今后每当 $w$ 由 $u$ 经\eqref{align:ch09-2-eq23}、\eqref{align:ch09-2-eq24} 导出时，写 $A[u]=w$。由\eqref{align:ch09-2-eq22} 与\eqref{align:ch09-2-eq25}，有估计
\begin{align}
\|A[u]\|_{H^2(U)}\leqslant C(\|u\|_{H_0^1(U)}+1).\label{align:ch09-2-eq26}
\end{align}
2. 现断言 $A:H_0^1(U)\to H_0^1(U)$ 连续且紧。事实上，若
\begin{align}
u_k\to u\quad\text{于 }H_0^1(U),\label{align:ch09-2-eq27}
\end{align}
则由估计\eqref{align:ch09-2-eq26}，
\begin{align}
\sup_k\|w_k\|_{H^2(U)}<\infty,\label{align:ch09-2-eq28}
\end{align}
其中 $w_k=A[u_k]$（$k=1,\cdots$）。故存在子列 $\{w_{k_j}\}_{j=1}^\infty$ 及函数 $w\in H_0^1(U)$ 使
\begin{align}
w_{k_j}\to w\quad\text{于 }H_0^1(U).\label{align:ch09-2-eq29}
\end{align}
现
\begin{align*}
\int_U Dw_{k_j}\cdot Dv+\mu w_{k_j}v\,\mathrm{d}x=-\int_U b(Du_{k_j})v\,\mathrm{d}x
\end{align*}
对每个 $v\in H_0^1(U)$。于是用\eqref{align:ch09-2-eq22}、\eqref{align:ch09-2-eq27} 与\eqref{align:ch09-2-eq29}，得
\begin{align*}
\int_U Dw\cdot Dv+\mu wv\,\mathrm{d}x=-\int_U b(Du)v\,\mathrm{d}x
\end{align*}
对每个 $v\in H_0^1(U)$。故 $w=A[u]$。
于是\eqref{align:ch09-2-eq27} 蕴涵 $A[u_k]\to A[u]$ 于 $H_0^1(U)$，从而 $A$ 连续。类似论证表明 $A$ 紧：若 $\{u_k\}_{k=1}^\infty$ 在 $H_0^1(U)$ 中有界，估计\eqref{align:ch09-2-eq22} 断言 $\{A[u_k]\}_{k=1}^\infty$ 在 $H^2(U)$ 中有界，故在 $H_0^1(U)$ 中有强收敛子列。
3. 最后须证明：若 $\mu$ 充分大，集合
\begin{align*}
\{u\in H_0^1(U)\mid u=\lambda A[u]\text{ 对某 }0\leqslant\lambda\leqslant 1\}
\end{align*}
在 $H_0^1(U)$ 中有界。设 $u\in H_0^1(U)$，$u=\lambda A[u]$，某 $0<\lambda\leqslant 1$。则 $\frac u\lambda=A[u]$；换言之，$u\in H^2(U)\cap H_0^1(U)$，且
\begin{align*}
-\Delta u+\mu u=-\lambda b(Du)\quad\text{a.e. 在 }U\text{ 中}.
\end{align*}
以 $u$ 乘此恒等式并在 $U$ 上积分，算得
\begin{align*}
\int_U|Du|^2+\mu|u|^2\,\mathrm{d}x=-\int_U\lambda b(Du)u\,\mathrm{d}x\leqslant\int_U C(|Du|+1)|u|\,\mathrm{d}x\leqslant\frac12\int_U|Du|^2\,\mathrm{d}x+C\int_U|u|^2+1\,\mathrm{d}x.
\end{align*}
故若 $\mu>0$ 充分大，$\|u\|_{H_0^1(U)}\leqslant C$，某常数 $C$ 不依赖于 $0\leqslant\lambda\leqslant 1$。
4. 在空间 $X=H_0^1(U)$ 中应用 Schaefer 不动点定理，推得 $A$ 有不动点 $u\in H_0^1(U)\cap H^2(U)$，后者进而求解半线性 PDE \eqref{align:ch09-2-eq21}。
\end{proof}
\end{theorem}
\end{example}

\paragraph{警告。}构造性求解\eqref{align:ch09-2-eq21} 的一个看似合理的方案是选取某个 $u_0$，然后迭代求解线性边值问题
\begin{align*}
\begin{cases}
-\Delta u_{k+1}+\mu u_{k+1}=-b(Du_k)&\text{在 }U\text{ 中},\\
u_{k+1}=0&\text{在 }\partial U\text{ 上}
\end{cases}\qquad (k=0,1,\cdots).
\end{align*}
然而不能断言 $\{u_k\}_{k=0}^\infty$ 然后收敛到\eqref{align:ch09-2-eq21} 的解。Schauder 与 Schaefer 不动点定理并不断言任何序列收敛到不动点。（但见 \S9.3 后续的证明。）

用不动点定理解非线性椭圆 PDE 的更精细应用，见 Gilbarg--Trudinger [G-T]。

\end{document}
